\chapter{Associative-node no-go: почему color-six не является Krajewski node}

Трёхузловой exterior-square carrier дал правильный determinant и Hessian.
Теперь проводится более строгая проверка: могут ли его nodes быть modules
той же ассоциативной Pati--Salam algebra, которая действует на fermions.
Ответ отрицателен, и это уточняет, где именно должна жить найденная
конструкция.

\section{Локализация obstruction}

Determinant-stabilizing identity сохраняется, но её literal realization через
новый шестимерный fermion module невозможна. Группа \(SU(4)\) действует на
шестёрке, тогда как полная matrix algebra \(M_4(\mathbb C)\) не допускает на
ней линейного unital associative representation.

Значит:
\begin{enumerate}
 \item формулы determinant и положительного Hessian остаются верны;
 \item literal node \((1_R,6_4)\) в строгом Krajewski graph запрещён;
 \item шестёрка должна появляться как scalar curvature channel внутри
 существующего Hilbert space, где она уже присутствует в Majorana block.
\end{enumerate}

\section{Группа не равна associative algebra}

Для \(g\in GL(4,\mathbb C)\) внешний квадрат задаёт корректное group
representation
\begin{equation}
 \Lambda^2(gh)=\Lambda^2(g)\Lambda^2(h).
\end{equation}
Его производная задаёт корректное Lie-algebra representation
\begin{equation}
 d\Lambda^2([X,Y])
 =[d\Lambda^2(X),d\Lambda^2(Y)].
\end{equation}
Именно поэтому предыдущий wedge edge является точно gauge-equivariant.

Но representation ассоциативной algebra должна быть complex-linear,
unital и multiplicative. Карта \(A\mapsto\Lambda^2A\) квадратична:
\begin{equation}
 \Lambda^2(A+B)\ne\Lambda^2A+\Lambda^2B,
 \qquad
 \Lambda^2(\lambda I_4)=\lambda^2I_6.
\end{equation}
Производная \(d\Lambda^2\) линейна, но
\begin{equation}
 d\Lambda^2(I_4)=2I_6,
 \qquad
 d\Lambda^2(AB)\ne d\Lambda^2(A)d\Lambda^2(B).
\end{equation}
Следовательно, ни одна из двух естественных exterior-square maps не является
unital associative representation \(M_4(\mathbb C)\to M_6(\mathbb C)\).

На двухстах random matrix pairs machine audit даёт ошибки порядка
\(10^{-14}\) для group multiplicativity и Lie commutators, но типичный
additivity/associative-product defect порядка \(10^1\). Scalar test уже
точно даёт
\begin{equation}
 \|\Lambda^2(2I_4)-2I_6\|_F=2\sqrt6,
 \qquad
 \|d\Lambda^2(I_4)-I_6\|_F=\sqrt6.
\end{equation}

\section{Artin--Wedderburn module obstruction}

Project Pati--Salam algebra равна
\begin{equation}
 \mathcal A_F^{PS}
 =\mathbb H_R\oplus\mathbb H_L\oplus M_4(\mathbb C).
\end{equation}
Её irreducible complex left modules имеют dimensions
\begin{equation}
 2,\qquad2,\qquad4.
\end{equation}
В частности, любой nontrivial \(M_4(\mathbb C)\)-module является direct sum
копий fundamental \(\mathbb C^4\), поэтому его dimension кратна четырём.
Irreducible color-six module и unital one-dimensional module отсутствуют.
Следовательно, оба крайних nodes chain
\begin{equation}
 \mathbb C\longrightarrow(2,4)\longrightarrow(1,6)
\end{equation}
не являются nodes строгого finite triple над неизменённой
\(\mathcal A_F^{PS}\).

Добавление отдельного \(M_6(\mathbb C)\)-summand формально разрешило бы
six-dimensional module, но одновременно добавило бы новый \(U(6)\) gauge
factor и тем самым изменило бы модель, которую требовалось вывести.

\section{Почему wedge-map всё равно не случайна}

На уровне gauge representations
\begin{equation}
 \operatorname{Hom}((2,4),(1,6))
 \cong(2,4)\oplus(2,\overline{20}).
\end{equation}
Wedge polarization выбирает единственную \((2,4)\)-копию. Machine map
\begin{equation}
 \Delta\longmapsto B_\Delta
\end{equation}
имеет complex rank восемь и проходит exact
\(SU(2)_R\times SU(4)\)-equivariance test. Поэтому проблема не в
representation content scalar field, а только в ошибочном превращении
scalar representation в fermionic algebra module.

\section{Правильное место для шестёрки}

В уже построенном finite block
\begin{equation}
 \mathcal H_F=V_R\oplus V_L\oplus V_R^c\oplus V_L^c
\end{equation}
color-six присутствует внутри
\begin{equation}
 \operatorname{Sym}^2(2_R\otimes4_4)
 =(3_R,10_4)\oplus(1_R,6_4),
\end{equation}
то есть как scalar Majorana channel, а не как Hilbert node. Это указывает
на корректный carrier: represented degree-two differential calculus
\begin{equation}
 \Omega^2_{D_F}(\mathcal A_F)
 =\pi_D(\Omega^2\mathcal A_F)/\pi_D(J^2)
\end{equation}
или соответствующая curvature superconnection. Два length-two paths могут
породить wedge component в \(D_F^2\), а junk quotient должен решить, остаётся
ли нужная antisymmetric projection в physical curvature.

\begin{nogocriterion}[Associative-node no-go]
Exterior-square chain нельзя объявлять Krajewski finite triple над
\(\mathbb H_R\oplus\mathbb H_L\oplus M_4(\mathbb C)\): representations
\(1\) и \(6\) не являются допустимыми endpoint modules этой associative
algebra. Coarse order-one label test был необходим, но недостаточен.
\end{nogocriterion}

\begin{proposition}[Сохранённый curvature route]
Exterior-square trace identity остаётся точным gauge-equivariant указателем
на нужный determinant selector. Допустимый следующий gate должен не добавлять
новые fermion nodes, а вычислить represented universal two-forms и junk ideal
на существующем 32-dimensional Pati--Salam module и проверить, выделяет ли
physical curvature канал \((1_R,6_4)\).
\end{proposition}

\section{Литература и воспроизводимость}

Различие между vertices finite triple и scalar links сверено с
T.~Krajewski, \emph{Classification of Finite Spectral Triples},
arXiv:hep-th/9701081. Ограничения scalar potential и роль diagram paths
сверены с T.~Krajewski, arXiv:hep-th/9803199. Generalized inner fluctuations
сверены с A.~H.~Chamseddine and W.~D.~van Suijlekom, arXiv:1304.7583, а
Pati--Salam scalar channels --- с A.~H.~Chamseddine, A.~Connes and
W.~D.~van Suijlekom, arXiv:1304.8050.
Переход от universal forms к represented calculus modulo junk является
стандартной частью spectral-triple differential calculus.

Machine audit сохранён в
\path{s2t_v4_pati_salam_associative_node_no_go.py} и
\path{s2t_v4_pati_salam_associative_node_no_go_results.json}.